\section*{الفصل \ref{ch:b1:newton-dynamics} --- قوانين نيوتن في الديناميك}

\begin{solution}{exo:b1:newton-dynamics:1}
قراءة الميزان $N = m(g + a)$ حيث $a$ التسارع نحو الأعلى: فإلى أعلى،
$70 \times 11.81 = \qty{827}{N}$ (أي ``\qty{84}{kg}'')؛ وإلى أسفل، $70 \times 7.81 =
\qty{547}{N}$ (أي ``\qty{56}{kg}'')؛ وبسرعة ثابتة، \qty{687}{N}
(أي \qty{70}{kg})؛ وفي السقوط الحرّ، $0$.
\end{solution}

\begin{solution}{exo:b1:newton-dynamics:2}
التسارع $a$ نفسه لكليهما (لأن الخيط غير قابل للاستطالة): $m_1a = T$ و
$m_2a = m_2g - T$، ومنه $a = m_2g/(m_1 + m_2) = \qty{3.3}{m/s^2}$ و
$T = m_1a = \qty{6.5}{N}$.
\end{solution}

\begin{solution}{exo:b1:newton-dynamics:3}
$a = g(\sin\alpha - \mu_d\cos\alpha) = 9.81(0.50 - 0.17) = \qty{3.2}{m/s^2}$؛
و$v = \sqrt{2 \times 3.2 \times 5.0} = \qty{5.7}{m/s}$. ولا أثر \omterm{def:b1:newton-dynamics:momentum}{للكتلة}.
\end{solution}

\begin{solution}{exo:b1:newton-dynamics:4}
$k\Delta\ell = mg$: ومنه $\Delta\ell = 0.50 \times 9.81/200 = \qty{2.5}{cm}$.
ومع $x$ الانتقالَ عن التوازن يكون $m\ddot x = -kx$ (إذ يتلاشى الوزن
مع الاستطالة السكونية): ومنه $\omega_0 = \sqrt{k/m} = \qty{20}{rad/s}$ و
$T = \qty{0.31}{s}$.
\end{solution}

\begin{solution}{exo:b1:newton-dynamics:5}
$\mu_s = \tan 22^\circ = 0.40$. وفي المنعطف: $v \leq \sqrt{\mu_sgR}$: أي $\sqrt{0.8
\times 9.81 \times 50} = \qty{20}{m/s}$ (أي \qty{71}{km/h})؛ وعلى المبتلّ،
\qty{14}{m/s} (أي \qty{50}{km/h}).
\end{solution}

\begin{solution}{exo:b1:newton-dynamics:6}
المظلّي: $v_\infty = \sqrt{2 \times 785/(1.2 \times 0.70)} = \qty{43}{m/s}$؛
و$\mathrm{Re} = \rho vL/\eta \approx 1.2 \times 43 \times 0.5/1.8\times10^{-5}
\approx 10^6$: أي تربيعي. والقطرة: $m = \tfrac43\pi r^3\rho_w = \qty{4.2e-12}{kg}$ و
$v_\infty = mg/6\pi\eta r = 4.1\times10^{-11}/3.4\times10^{-9} = \qty{1.2}{cm/s}$؛
و$\mathrm{Re} = 1.2 \times 0.012 \times 2\times10^{-5}/1.8\times10^{-5} = 0.02$:
أي خطي.
\end{solution}

\begin{solution}{exo:b1:newton-dynamics:7}
عند الأسفل: $v^2 = 2g\ell(1 - \cos 60^\circ) = g\ell$؛ ومنه $T = mg + mv^2/\ell =
2mg$. وعند نقطة الإفلات: $v = 0$ و $T = mg\cos\theta_0 = 0.5mg$.
\end{solution}

\begin{solution}{exo:b1:newton-dynamics:8}
$v = \sqrt{gR\tan\beta} = \sqrt{9.81 \times 200 \times 0.268} = \qty{23}{m/s}$
(أي \qty{82}{km/h}). وعند سرعة أدنى تميل السيارة إلى الانزلاق نازلةً على الميل، فيجب أن يشير
الاحتكاك صاعدًا؛ وعند سرعة أعلى يشير نازلًا ليضيف
\omterm{def:b1:newton-dynamics:momentum}{قوة} مركزية.
\end{solution}

\begin{solution}{exo:b1:newton-dynamics:9}
$T\cos\alpha = mg$ و $T\sin\alpha = m\omega^2\ell\sin\alpha$: ومنه $\omega =
\sqrt{g/\ell\cos\alpha} = \sqrt{9.81/0.866} = \qty{3.4}{rad/s}$، والدور
\qty{1.9}{s}؛ و$T = mg/\cos\alpha = 1.15mg$.
\end{solution}

\begin{solution}{exo:b1:newton-dynamics:10}
$a = 27.8/8.0 = \qty{3.5}{m/s^2}$ و $F = \qty{3.5}{kN}$؛ ويعطي $F \leq \mu_smg$
أن $\mu_s \geq a/g = 0.35$. وعلى طريق مبتلّ يحدّ $\mu_s \approx 0.4$
التسارعَ قرب $0.4g$ مهما يكن المحرك: فالتماسك لا الاستطاعة.
\end{solution}

\begin{solution}{exo:b1:newton-dynamics:11}
$v = v_\infty(1 - \eu^{-t/\tau})$ و $v_\infty = g\tau = \qty{4.9}{m/s}$؛
و$z = \int v = v_\infty[t - \tau(1 - \eu^{-t/\tau})]$. وعند \qty{2.0}{s}:
$4.9 \times (2.0 - 0.5 \times 0.98) = \qty{7.4}{m}$، في مقابل $\tfrac12 gt^2
= \qty{20}{m}$ في الخلاء.
\end{solution}

\begin{solution}{exo:b1:newton-dynamics:12}
الإسقاط القطري (نحو المركز): $N - mg\cos\theta = mv^2/R$، ومنه
$N = m(v^2/R + g\cos\theta)$. وعند القمة (حيث $\theta = \pi$): $N = m(v^2/R - g)
\geq 0$ إذا وفقط إذا كان $v \geq \sqrt{gR}$. ومع قانون قيمة السرعة، $v_{\text{القمة}}^2 =
v_0^2 - 4gR \geq gR$: أي $v_0 \geq \sqrt{5gR}$؛ وعندئذٍ يكون عند الأسفل $N =
m(5g + g) = 6mg$.
\end{solution}

\begin{solution}{pb:b1:newton-dynamics:1}
\textbf{1.} $m\dot v = mg$: ومنه $v = gt$ و $z = \tfrac12 gt^2$؛ وعند \qty{10}{s}:
\qty{98}{m/s} و\qty{490}{m}.

\textbf{2.} $mg = \qty{785}{N}$.

\textbf{3.} على الأرض تدفع الأرضية إلى أعلى \omterm{def:b1:newton-dynamics:momentum}{بالقوة} $N = mg$ --- وذلك الدفع
هو ما يُشعر به وزنًا. أما في السقوط الحرّ فلا \omterm{def:b1:newton-dynamics:momentum}{قوة} تماسّ:
إذ تؤثر الجاذبية في كل جزء تأثيرًا واحدًا ولا يضغط شيء.

\textbf{4.} نعم إلى حدود جزء من المئة على مدى دقيقة؛ والمهمَل:
دوران الأرض (أي الحدّان الطارد المركزي وكوريوليس،
\cref{ch:b1:non-inertial-frames}) والرياح.

\textbf{5.} $F_d = \tfrac12\rho C_xSv^2 = 0.48\,v^2$ (بوحدات النظام الدولي).

\textbf{6.} $m\,\dd v/\dd t = mg - 0.48v^2$.

\textbf{7.} $v_\infty = \sqrt{mg/0.48} = \sqrt{1635} = \qty{40}{m/s} =
\qty{146}{km/h}$.

\textbf{8.} نقسم على $mg$: $\dd u/\dd t\,(v_\infty/g) = 1 - u^2$؛ و
$\tau = 40.4/9.81 = \qty{4.1}{s}$.

\textbf{9.} $\int_0^u\dd u'/(1 - u'^2) = \operatorname{artanh}u = t/\tau$،
ومنه $u = \tanh(t/\tau)$.

\textbf{10.} $t = \tau\operatorname{artanh}(0.9) = 1.47\tau = \qty{6.1}{s}$؛
ومن أجل $0.99$: $2.65\tau = \qty{11}{s}$.

\textbf{11.} $z = \int v_\infty\tanh(t/\tau)\dd t = v_\infty\tau\ln\cosh(t/\tau)$؛
وعند \qty{10}{s}: $166 \times \ln\cosh(2.43) = 166 \times 1.74 = \qty{290}{m}$،
في مقابل \qty{490}{m}.

\textbf{12.} $\mathrm{Re} = 1.2 \times 40 \times 0.5/1.8\times10^{-5} \approx
\num{1.3e6}$: أي أعلى بكثير من $1000$، فتربيعي.

\textbf{13.} $v_\infty' = \sqrt{785/(0.6 \times 25)} = \sqrt{52} = \qty{7.2}{m/s}$.

\textbf{14.} السحب $= 15 \times 1600 = \qty{24}{kN}$؛ و$a = (24000 - 785)/80
= \qty{290}{m/s^2} \approx 30g$: فينكسر الحزام، أو ينكسر القافز.

\textbf{15.} متوسط $a \approx (40 - 7)/3 = \qty{11}{m/s^2} \approx 1.1g$؛
\omterm{def:b1:newton-dynamics:momentum}{وقوة} الحزام $\approx m(g + a) \approx \qty{1.7}{kN}$.

\textbf{16.} $\tau' = 7.2/9.81 = \qty{0.73}{s}$: أي إن النزول يصير مستقرًّا خلال
بضع ثوانٍ من الانفتاح الكامل.

\textbf{17.} $\vect a = \vect 0$: \omterm{def:b1:newton-dynamics:momentum}{فقوة} الحزام $= mg = \qty{785}{N}$.

\textbf{18.} يعطي التوازن $mg = \tfrac12\rho C_xSv^2$ أن $v \propto (C_xS)^{-1/2}$:
فنصف المساحة يعطي $\sqrt2$ ضعف السرعة (أي \qty{10}{m/s}).

\textbf{19.} $v \propto \sqrt m$: ومنه $7.2 \times \sqrt{100/80} = \qty{8.1}{m/s}$.

\textbf{20.} $h = v^2/2g = 52/19.6 = \qty{2.6}{m}$: أي قفزة من الطابق
الأول.

\textbf{21.} $a = v^2/2d$: أي \qty{52}{m/s^2} ($5.3g$) على \qty{0.5}{m}،
\omterm{def:b1:newton-dynamics:momentum}{والقوة} $m(g + a) \approx \qty{5}{kN}$؛ و\qty{17}{m/s^2} ($1.8g$) على
\qty{1.5}{m}، أي \qty{2.2}{kN}. فالتدحرج.

\textbf{22.} سحب أكبر (مع بعض الرفع) يتجاوز الوزن برهةً:
فيتباطأ القافز ويلامس الأرض دون $v_\infty'$.

\textbf{23.} $m = \tfrac43\pi r^3\rho_w = \qty{4.2e-6}{kg}$ و $S = \pi r^2 =
\qty{3.1e-6}{m^2}$: ومنه $v_\infty = \sqrt{2mg/\rho C_xS} = \sqrt{8.2\times10^{-5}/
1.9\times10^{-6}} = \qty{6.6}{m/s}$؛ و$\mathrm{Re} = 1.2 \times 6.6 \times
2\times10^{-3}/1.8\times10^{-5} \approx 900$: أي تربيعي (تقريبًا).

\textbf{24.} $v_\infty = \qty{1.2}{cm/s}$ (\cref{exo:b1:newton-dynamics:6})؛
و$\mathrm{Re} \approx 0.02 \ll 1$: فينطبق ستوكس، على نحو متسق.

\textbf{25.} $m\dot{\vect v} = m\vect g + \vect F_{\text{السحب}}$؛ والسحب خطي
في $v$ للأشياء الصغيرة البطيئة، وتربيعي للكبيرة السريعة. و
الأعداد: \qty{40}{m/s} لجسم ساقط، و\qty{7}{m/s} تحت مظلّة
(أو لقطرة مطر)، و\qty{1}{cm/s} للضباب --- وهو القانون نفسه ممتدًّا على أربع
رتب من رتب القدر.
\end{solution}
